\section{相关工作}
\label{sec:related}

\subsection{目标检测中的知识蒸馏}

知识蒸馏由\cite{hinton2015distilling}提出，通过匹配软化的输出分布将大型教师网络的知识转移到小型学生网络。对于目标检测，挑战超出了分类，还包括定位和特征学习。

\cite{chen2017learning}开创了检测领域的KD研究，提出同时蒸馏分类和边界框回归输出。后续工作探索了各种知识来源：FitNets\cite{romero2014fitnets}关注中间特征对齐；注意力迁移方法\cite{zagoruyko2016paying}使用注意力图作为监督；定位感知蒸馏\cite{dai2021general}专门针对边界框质量。

近期检测KD框架通常结合多种蒸馏分支，集成特征和注意力蒸馏以实现空间一致性，而基于关系的蒸馏则捕获结构信息。然而，这些方法通常在干净数据集（COCO、ImageNet）上验证，在恶劣条件下的行为研究不足。

\subsection{域偏移与恶劣条件下的知识蒸馏}

KD对域偏移的鲁棒性已引起越来越多的关注。\cite{sun2016deep}分析跨域的特征对齐，而小样本适应方法则探索在有限目标数据下的跨域知识转移\cite{mirzadeh2020improved}。

恶劣天气条件构成特别有挑战性的域偏移场景。\cite{guo2021meets}综述了KD与图像复原的交叉研究，指出雾和霾等退化会降低教师和学生模型的性能。\cite{du2022learning}提出通过对抗训练学习天气不变表示，但没有分离具体的失效机制。

若干工作已针对雾天目标检测。\cite{li2018benchmarking}建立了去雾和可见度复原的基准。\cite{sakaridis2018semantic}提供了Foggy Cityscapes数据集，证明合成雾天下性能显著下降。这些工作虽承认标准检测器在雾天 struggling，但未分析知识转移具体如何失效。

\subsection{机制导向的知识蒸馏分析}

超越提出新方法，日益增多的工作试图理解KD\emph{为什么}有效或失效。\cite{phuong2019functional}对蒸馏学习的函数进行理论分析，而\cite{cho2019efficacy}检查教师知识的哪些部分实际发生转移。

实证机制研究通常聚焦教师-学生关系。\cite{mirzadeh2020improved}调查教师和学生能力差距，发现不完美的教师仍可提供有用的监督。\cite{gou2021knowledge}综述各种知识类型及其转移机制。

对于失效分析，若干假设在文献中反复出现。\textbf{分布错配}假设域偏移破坏了KD依赖的教师输出分布\cite{wang2020cross}。\textbf{不确定性放大}假设教师在退化输入下预测变得不可靠，在学生训练中放大误差\cite{kendall2017uncertainties}。\textbf{表征不对齐}假设特征空间跨域发散，阻碍基于特征的蒸馏\cite{sun2016deep}。

然而，这些机制很少在统一框架内相互测试。我们的工作通过在控制性实验设计中实证评估多种假设来填补这一空白，确定哪些机制得到观测行为的直接支持。
